% Generated draft: requires mathematical and source review.
\documentclass[11pt]{article}
\usepackage[margin=25mm]{geometry}
\usepackage{fontspec}
\setmainfont{Latin Modern Roman}
\usepackage{amsmath,amssymb,mathtools,unicode-math}
\setmathfont{Latin Modern Math}
\usepackage{enumitem,needspace,xurl,hyperref}
\setlength{\parindent}{0pt}
\setlength{\parskip}{5pt}
\setlength{\emergencystretch}{3em}
\newcommand{\N}{\mathbb{N}}
\newcommand{\Z}{\mathbb{Z}}
\newcommand{\Q}{\mathbb{Q}}
\newcommand{\R}{\mathbb{R}}
\newcommand{\C}{\mathbb{C}}
\newcommand{\F}{\mathbb{F}}
\newcommand{\A}{\mathbb{A}}
\newcommand{\PP}{\mathbb{P}}
\newcommand{\E}{\mathbb{E}}
\newcommand{\eps}{\varepsilon}
\newcommand{\dd}{\,\mathrm{d}}
\begin{document}
\section*{UnsolvedMath ID: 30002349}
\textbf{Source problem code:} catalogue entry
\subsection{Definitions and mathematical statement}
Let \(p\) be a prime, let \(k\) be an algebraically closed field of characteristic \(p\), and let \(k[[t]]\) and \(k((t))\) denote the formal power series ring and its fraction field. Let \(L/k((t))\) be a Galois extension with group \(G\), and put
\[
A=\{a\in L:a\text{ is integral over }k[[t]]\}.
\]
The source asks whether this local extension admits a characteristic-zero lift.

For a precise working formulation, assume that \(L/k((t))\) is finite, so that \(G\) is finite. This is the usual finite local-extension interpretation; finiteness is not explicit in the supplied statement. Under this interpretation, ask whether, for every such \(p,k,L\), there exist a complete discrete valuation ring \(R\) of characteristic zero, its maximal ideal \(\mathfrak m_R\), an identification \(R/\mathfrak m_R\cong k\), and a finite flat integral \(R[[T]]\)-algebra \(B\) equipped with an action of \(G\) by \(R[[T]]\)-algebra automorphisms, such that
\[
B^G=R[[T]],\qquad
\operatorname{Frac}(B)/\operatorname{Frac}(R[[T]])
\text{ is Galois with group }G,
\]
and there is an isomorphism
\[
\phi:B/\mathfrak m_RB\xrightarrow{\ \sim\ }A
\]
of \(k[[t]]\)-algebras, where the base identification
\(R[[T]]/\mathfrak m_RR[[T]]\cong k[[t]]\) sends \(T\) to \(t\). The induced generic Galois action is required to be the specified action of \(G\) on \(B\).

This spells out a finite-algebra interpretation of the undefined phrase ``\(G\)-Galois lift.'' If the original Galois action itself must lift, require in addition
\[
\phi(g\cdot\overline b)=g\cdot\phi(\overline b)
\qquad(g\in G,\ \overline b\in B/\mathfrak m_RB),
\]
using the action on \(A\) induced from \(L\). The supplied wording explicitly requests a \(k[[t]]\)-algebra isomorphism but does not explicitly state this equivariance requirement. These interpretations require source clarification.

The supplied background also mentions an additional refinement controlling \(R\), but gives no condition expressing that refinement; it cannot be formulated precisely from this record. The supplied open label and literature assessment are dated curation, not mathematical conclusions established by this draft.
\subsection{Short English statement}
Given a local Galois extension in characteristic \(p\), the question asks whether it arises by reducing a Galois extension over a complete discrete valuation ring of characteristic zero with residue field \(k\). The finite-extension convention, compatibility with the specified \(G\)-action, and an unprovided refinement controlling the lifting ring require clarification.
\subsection{Sources}
\begin{itemize}
\item Supplied local UnsolvedMath registry record: permanent ID \(30002349\), source problem code \texttt{catalogue entry}, title ``Characteristic-Zero Lifts of Local Galois Extensions,'' in the Oberwolfach Reports Open Problems collection. No individual proposer is supplied; the proposed year is recorded as 2013.
\item Supplied source citation: \emph{The Arithmetic of Fields (2014)}, Oberwolfach Reports, DOI \texttt{10.4171/owr/2013/30}, pp.\ 1771--1818. \url{https://doi.org/10.4171/owr/2013/30}. The background describes these reports as records of mathematical workshops and their open-problem discussions.
\item Supplied literature reference: Pop, \emph{The Oort conjecture on lifting covers of curves}. \url{https://annals.math.princeton.edu/2014/180-1/p06}.
\item Supplied literature reference: Obus--Wewers, \emph{Cyclic extensions and the local lifting problem}. \url{https://annals.math.princeton.edu/2014/180-1/p05}.
\item Supplied literature reference: \emph{The Arithmetic of Fields}, Oberwolfach Reports. \url{https://doi.org/10.4171/owr/2013/30}.
\item Literature assessment dated exactly \texttt{2026-08-21}: the supplied curation labels the problem open, states that the cyclic-group local lifting problem is solved by the Oort conjecture theorem, and describes the arbitrary-group question and an additional ring-control refinement as unresolved. This is supplied dated curation derived from the OWR material, not independently verified here; the record itself requests independent verification of time-sensitive status and cited claims.
\end{itemize}
\end{document}
